%%
%% spring-lecture09.tex
%% 
%% Made by Alex Nelson <pqnelson@gmail.com>
%% Login   <alex@lisp>
%% 
%% Started on  2026-04-28T09:40:48-0700
%% Last update 2026-04-28T09:40:48-0700
%% 

\lecture

\begin{lemma}\label{lemma:3.21}
Let $\FF$ be an algebraically closed field, let $D$ be a
finite-dimensional division algebra over $\FF$. Then $D=\FF$.
\end{lemma}

\begin{corollary}\label{cor:3.22}
Let $\FF$ be an algebraically closed field, let $A$ be a
finite-dimensional $\FF$-algebra.
Then $A$ is semisimple if and only if
$A\iso\prod^{r}_{i=1}M_{n_{i}}(\FF)$ where $1\leq n_{1}\leq\dots\leq n_{r}$.
Moreover, $A$ is simple if and only if $A\iso M_{n}(\FF)$.
\end{corollary}

\begin{theorem}[Maschke]\label{thm:3.23}
Let $G$ be a finite group, let $\FF$ be a field. Then the group
algebra $\FF[G]$ is semisimple if and only if $\Char(\FF)\ndivides|G|$.
\end{theorem}

\begin{proof}
\backwardproof\ Assume $\Char(\FF)\ndivides n$ where $n=|G|$.
(Crucially, this means $n$ is a unit in $\FF[G]$.)
By Proposition~\ref{prop:2.15}, it suffices to show
that the lattice of submodules $S(\FF[G])$ is complemented. For a
right ideal $M$ of $\FF[G]$, we want to find its complement---if there
exists an $\FF[G]$-module morphism
\begin{equation}
\varphi\colon\FF[G]\to M
\end{equation}
such that $\varphi(u)=u$ for all $u\in M$, then $N=\ker(\varphi)=(1-\varphi)M$
is a right ideal of $\FF[G]$ such that
\begin{equation}
M\oplus N=\FF[G].
\end{equation}
We want to construct this. We start with the $\FF$-linear map
\begin{equation}
\pi\colon\FF[G]\to m
\end{equation}
such that $\pi(u)=u$ for all $u\in M$. We define $\varphi$ by
``averaging'' $\pi$ over $G$,
\begin{equation}
\varphi(u):=\left(\sum_{x\in G}\pi(ux)x^{-1}\right)n^{-1}
\end{equation}
for all $u\in\FF[G]$ and $n=|G|$. (Note: $u\in\FF[G]$ has $u^{-1}$
exist by hypothesis $\Char(\FF)\ndivides|G|$) Then $\varphi$ satisfies
\begin{equation}
\varphi(uy)=\varphi(u)y
\end{equation}
for all $u\in\FF[G]$ and $y\in G$ [i.e., $\varphi$ satisfies a sort of
  linearity], and $\varphi(u)=u$ for all $u\in M$. Hence the claim.

\forwardproof\ (We will prove the contrapositive.)
Assume that $\Char(\FF)\divides n$ where $n=|G|$. Then for any
$u\in\FF[G]$ we see that $nu=0$. In particular, the nonzero element
$e\in\FF[G]$ defined by
\begin{equation}
e:=\sum_{g\in G}g
\end{equation}
we see that
\begin{equation}
ey=e
\end{equation}
for all $y\in G$. Therefore
\begin{equation}
e^{2}=0
\end{equation}
[Pf: $e^{2}=\sum_{g\in G}eg=\sum_{g\in G}e=ne=0$.]
Moreover, the right ideal $N$ generated by $e$ satisfies
\begin{equation}
N=e\FF,
\end{equation}
and thus
\begin{equation}
N^{2}=e^{2}\FF=0.
\end{equation}
By Lemma~\ref{lemma:3.7},
\begin{equation}
0\neq N\subset\Rad(\FF[G]),
\end{equation}
and by Proposition~\ref{prop:3.2}~(1) $\FF[G]$ is
non-semisimple. Hence the contrapositive.
\end{proof}

\begin{remark}
Maschke's theorem is one of the most profound theorems concerning the
representation of (finite) groups. See Curtis and Reiner~\cite{curtis1962representation}
(or any other decent book about representation theory) for more about
its consequences.
\end{remark}

\begin{corollary}\label{cor:3.24}
Let $G$ be a finite gorup, let $\FF$ be an algebraically-closed field
such that $\Char(F)\ndivides|G|$. Then $\FF[G]\iso\prod^{r}_{i=1}M_{n_{i}}(\FF)$
where $r$ is the number of conjugacy classes of $G$.
\end{corollary}

\begin{proof}
By Corollary~\ref{cor:3.22} and Theorem~\ref{thm:3.23}.
\end{proof}

\subsection{More about radicals}

\begin{fact}
$\Rad(A)=\bigcap\{N\in S(M)\mid M/N\mbox{ is simple}\}$.
\end{fact}

\begin{lemma}\label{lemma:4.1}
For [right] $A$-modules $M_{1}$ and $M_{2}$, and for any $A$-module
morphism $\phi\colon M_{1}\to M_{2}$, we have $\phi(\Rad(M_{1}))\subset\Rad(M_{2})$
and $\phi$ induces a morphism $M_{1}/\Rad(M_{1})\to M_{2}/\Rad(M_{2})$.
\end{lemma}

\begin{proof}
For any $N\subset M_{2}$ (i.e., $N\in S(M_{2})$),
$M_{1}/\phi^{-1}(N)\to M_{2}/N$ is injective [Pf: $\ker(\phi)=\phi^{-1}(0)\subset\phi^{-1}(N)$].
In particular, if the right-hand side is simple, then
$\phi^{-1}(N)=M_{1}$ or $M_{1}/\phi^{-1}(N)\iso M_{2}/N$ is simple. In
both cases, this $\phi^{-1}(N)\supset\Rad(M_{1})$, and this means
$\phi^{-1}(\Rad(M_{2}))\supset\Rad(M_{1})$. In other words, $\phi(\Rad(M_{1}))\subset\Rad(M_{2})$.
\end{proof}

\begin{proposition}\label{prop:4.2}
Let $A\neq0$ be some algebra. Then $\Rad(A)\in I(A)$ is a two-sided
ideal of $A$, and $\Rad(A)\in I(A)\setminus S(A)$ is a ``proper ideal''.
\end{proposition}

\begin{proof}
Check $\Rad(A)$ is a two-sided ideal: For any $x\in A$, the mapping
$(y\mapsto xy)\in\End_{A}(A)$ applied with Lemma~\ref{lemma:4.1}
implies $x(\Rad(A))\subset A$ [i.e., $A\cdot(\Rad(A))\subset A$,
hence it is a left ideal, hence 2-sided ideal].
  
Next we check $\Rad(A)\notin S(A)$. Since $1_{A}\in A$, by Zorn's Lemma
there exists a maximal proper 2-sided ideal $M\propersubset A$ (i.e.,
$A/M$ is simple) and thus $\Rad(A)\subset M\propersubset A$.
\end{proof}

\begin{corollary}\label{cor:4.3}
\begin{enumerate}
\item If $A$ is right Artinian, then $A/\Rad(A)$ is a semisimple algebra.
\item If $M$ is a right $A$-module, then $M(\Rad(A))\subset\Rad(M)$.
\end{enumerate}
\end{corollary}

\begin{proof}
\begin{enumerate}
\item By Lemma~\ref{lemma:2.22}, $A/\Rad(A)$ is also right
  $A$-Artinian (as an $A$-module) and by Lemma~\ref{lemma:2.28}~(3)
  $\Rad(A/\Rad(A))=0$. Then by Proposition~\ref{prop:3.2}~(1)
  $A/\Rad(A)$ is semisimple.
\item For any element $u\in M$, the mapping $(x\mapsto ux)\in\End_{A}(AM)$
  is a morphism. By Lemma~\ref{lemma:4.1},
  $u(\Rad(A))\subset\Rad(M)$. In other words, $M(\Rad(A))\subset M$.
\end{enumerate}
\end{proof}

\begin{lemma}\label{lemma:4.4}
Let $M$ be an $A$-module. For any $u\in M$, we have $u\in\Rad(M)$ if
and only if for every $N\subset M$ such that $uA+N=M$ we have $N=M$.
\end{lemma}

\begin{proof}
\backwardproof\ By contrapositive. If $u\notin\Rad(M)$, then there
exists an $N\subset M$ such that $M/N$ is simple and $u\notin N$. In
this case, $uA+N=M\neq N$.

\forwardproof\ By contrapositive. If there exists $N\subset M$ such
that $uA+N=M\neq N$ [i.e., $u\notin N$], then by Zorn's lemma there
exists a maximal submodule $P\subset M$ such that $N\subset P$ and
$u\notin P$. If there exists some submodule $Q$ such that
$P\propersubset Q$, then $u\in Q$ and thus $M=uA+N\subset Q$ (i.e., $M=Q$).
Therefore $M/P$ is simple and thus $\Rad(M)\subset P$ (since $u\notin P$
implies $u\notin\Rad(M)$).
\end{proof}

\begin{remark}
The radical of a module is analogous to the Frattini subgroup of a
group, and Nakayama's lemma is a variant of the usual characterization
of the Frattini subgroup.
\end{remark}

\begin{lemma}[Nakayama for modules]\label{lemma:4.5}
For any submodule $P\subset M$ such that
\begin{equation}\label{eq:NAK-for-modules}
\mbox{for any submodule }N\subset M\mbox{ such that }P+N=M\mbox{ we
  have }N=M.
\end{equation}
Then $P\subset\Rad(M)$.

Conversely, if $P\subset M$ and $P\subset\Rad(M)$, and either $P$ or
$M$ is finitely-generated as $A$-modules, then $P$ satisfies Equation~\eqref{eq:NAK-for-modules}.
\end{lemma}

\begin{lemma}[Nakayama for algebras]\label{lemma:4.6}
For any right ideal $P$ of $A$, the following are equivalent:
\begin{enumerate}
\item $P\subset\Rad(A)$
\item If $M$ is finitely-generated $A$ module and $N\subset M$ is a
  submodule such that $N+MP=M$, then $N=M$
\item $G=\{1+x\mid x\in P\}\subset A^{\times}$ is a subgroup of the
  group of units of $A$.
\end{enumerate}
\end{lemma}

\begin{corollary}[Nakayama's lemma]\label{cor:4.7}
Let $M$ be a finitely-generated $A$ module. If $M(\Rad(A))=M$, then $M=0$.
\end{corollary}

\begin{proof}
Take $N=0$ and $P=\Rad(A)$ in Lemma~\ref{lemma:4.6}~(2).
\end{proof}